\chapter{World Interaction}
\label{ch:world-interaction}
\index{world interaction}

\begin{tcolorbox}[colback=blue!5!white, colframe=blue!75!black, title=Using This Chapter,breakable]
This chapter gives you the tools to move through the world, interact with its people, and make your journey meaningful. For detailed descriptions of each region's culture, subgenre, and notable NPCs, see \textbf{Chapter \ref{ch:lore-compendium} (World Regions and Cultures)}. The two chapters work together: this one covers \emph{how} to travel and interact; the other tells you \emph{where} and \emph{with whom}.
\end{tcolorbox}

In \textbf{Fate's Edge}, the world is not a backdrop -- it is a partner in the conversation. Dikes groan under black rain in Viterra, clan horns answer across Acasia's ridgelines, Ecktoria's marble halls echo with careful words, and Kahfagia's pilots read storms by taste. Wherever you go, place, culture, and pressure push back.

\section{The Known World: A Geographic Overview}
\label{sec:known-world}
\index{geography}\index{regions}

The following summary orients the major regions relative to one another. (See Chapter \ref{ch:lore-compendium} for cultural details.)

\begin{description}
  \item[Ecktoria] \index{Ecktoria} Temperate heartland of the old Utaran Imperium. Borders Acasia to the northeast and Kahfagia to the west. Marble cities, imperial roads, codified law.
  \item[Acasia] \index{Acasia} Lowland of rivers and petty kingdoms, northeast of Ecktoria. Broken province of mercenary companies and border lords. Borders Vhasia to the east.
  \item[Vhasia] \index{Vhasia} Fractured kingdom on a southern peninsula. East of Acasia, west of Viterra. Courtly intrigue, martial tradition, and fractured sun symbolism.
  \item[Viterra] \index{Viterra} Shares the same peninsula with Vhasia, lying further east. Its eastern coast borders a dark inland body (the Blackwater). Hedge-law, river tolls, and the Last Kingdom.
  \item[Kahfagia] \index{Kahfagia} West of Ecktoria, facing the open ocean. Maritime empire of storm-flags, deep-sea pilots, and exotic imports.
  \item[Thepyrgos] \index{Thepyrgos} Occupies the southern part of the great peninsula opposite Vhasia/Viterra. Eastern coast overlooks the Astroegro Straits. City of a Thousand Stairs, center of learning.
  \item[Theona] \index{Theona} Large island off the eastern coast of Viterra. Three isles, causeways, marsh, and the ``No Ninth'' taboo.
  \item[Mistlands] \index{Mistlands} Western shore of the Aberderrian Sea (north of the Amaranthine). Bog, fog, and bell-lines; protectorate of the Aeler. Direwood lies to the west.
  \item[Linn] \index{Linn} Scattered around the Yloka River and the Aberderrian Sea. Longships reach the Blackwater and raid Theona and Viterra. Fjords, skerries, storm-oaths.
  \item[Violet Steppe] \index{Violet Steppe} North of the Amaranthine, west of the Aberderrian Sea. Vast grassland, home of the Ykrul nomads.
  \item[Valewood] \index{Valewood} East of the Blackwater. Ancient forest of phasing ruins, star-roads, and fae courts. Former Lethai empire.
  \item[Aelinnel (Aevrossa)] \index{Aelinnel} South of the Valewood. Gnomish homeland of stone, bough, and bright things. Mist paths, bell-mounds, tide-gates.
  \item[Aelaerem] \index{Aelaerem} South of Aevrossa. Halfling downs of gentle valleys and deep barrows. Homeland spoken of only in euphemisms (``The Silent Welcome'').
  \item[Zakov] \index{Zakov} Island in the Blackwater, near the mouth of the Yloka River. Pirate haven, reef passages, salt and serpent.
  \item[Oshiira] \index{Oshiira} Southern grasslands east of the Gulf of Oshiira. Grain confederacy, weigh-houses, and the Vermilion Corps.
  \item[Ashaan] \index{Ashaan} Eastern desert peninsula, south of the Crimson Valley. Veiled Throne, spirit-binding, and the Red Desert's creeping glass.
  \item[Sekogo] \index{Sekogo} Jungle realm south of the Crimson Valley and west of Ashaan. Green Council, honey-orchards, and the Ukwe (Speaking Tree).
  \item[Taharka] \index{Taharka} Highland kingdom of the Monsoon Crown. Coffee terraces, ridge-watchers, and passes that remember every army.
  \item[Fhara] \index{Fhara} Desert peninsula of oasis courts and water-clerks. Well-judges, coffee rituals, and the Sunken City of Qalat al-Madina.
  \item[Sidhi] \index{Sidhi} Diaspora people of the Brass Coast. Lighthouse courts, freedman's seals, and the green flame of memory.
  \item[Pereshi] \index{Pereshi} Highland fire-temple nation. Ash Prince's seal, fire-readers, and the ninth coal that never cools.
  \item[Kuvani] \index{Kuvani} Eastern steppe remount clans. Horse lineages, wind-songs, and the ghost herd of the salt pan.
  \item[Tulkani] \index{Tulkani} Itinerant storytellers and dream-carriers. Wagon circles, memory-cords, and the Dreamer who sleeps beneath the road.
  \item[Ngomebe] \index{Ngomebe} Walking cities of the southern plateau. Keystone chargers, roller-pits, and the Still Voice of abandoned stone.
  \item[Dhahara (Nine Kingdoms)] \index{Dhahara} Elephant kingdoms of the monsoon lowlands. Peacock Throne, river wars, and the ninth kingdom that was never whole.
\end{description}

\subsection*{The Ways Between}
\label{sec:ways-between}
\index{Ways Between}\index{Obishaal}

Beyond the physical map lie the Ways Between -- the dream-roads and spirit-paths of Obishaal, realm of dreams and death. Travel here is measured not in miles but in choices. A wrong turn might lead to yesterday. A right turn might lead to a version of yourself who made a different decision and is still alive to regret it.

\paragraph{Entering the Ways:} Thresholds exist in thin places -- crossroads at midnight, bridges that appear only at low tide, caves behind waterfalls, or simply a mist that settles too thick and stays too long. A character may also enter through ritual, dream, or near-death experience.

\paragraph{Navigation in Obishaal:} Use the standard travel procedures, but the GM may replace the Spades (place) draw with a \textbf{Virtue Test}. The road judges. Acts of honesty, courage, mercy, or self-sacrifice advance a personal \textbf{Virtue Timer [3]}; acts of greed, cowardice, cruelty, or despair advance a \textbf{Vice Timer [3]}. When one fills, the character ascends (Virtue) to a higher layer or falls (Vice) to a lower one.

\paragraph{The Thirteenth Stone:} At the deepest layer, pressed into the silt of the oldest dream, lies the Thirteenth Stone -- warm, pulsing, and whispering. Here, Ikasha, She Who Sleeps, dreams of freedom. Her voice offers bargains. Her freedom would remake the world.

\section{Game Structure and Time}
\label{sec:game-structure}
\index{Time}\index{Game structure}

Understanding how time works in Fate's Edge helps you navigate both the mechanical and narrative flow of play.

\subsection*{Basic Units}
\label{sec:time-units}
\begin{description}
\item[\textbf{Scene}] \index{Scene} The basic unit of narrative play, covering a specific situation or conflict (Some Time to Significant Time). Resolves a particular question or challenge.
\item[\textbf{Player Turn (Beat)}] \index{Turn} An individual player's action within a scene: Declare action $\rightarrow$ GM sets position $\rightarrow$ roll $\rightarrow$ resolve outcome $\rightarrow$ manage consequences.
\item[\textbf{Round}] \index{Round} Simultaneous or near-simultaneous actions within a scene (primarily for combat), representing a few seconds of real time.
\item[\textbf{Session}] \index{Session} One complete game session (typically 3--6 hours), containing 2--4 major scenes and resolving significant narrative progress.
\item[\textbf{Downtime}] \index{Downtime} The narrative time between scenes, used for recovery, advancement, and off-screen activities. Measured in days, weeks, or months depending on fiction.
\item[\textbf{Campaign}] \index{Campaign} Entire story arc (6--20+ sessions) with major character development and lasting consequences.
\end{description}

\section{Movement and Positioning}
\label{sec:movement-positioning}
\index{Movement}\index{Position}

Space is tracked with \textbf{range bands} and \textbf{Position} (see Chapter \ref{ch:gm-core} for full rules).

\subsection*{Range Bands}
\index{Range}
\begin{description}
\item[\textbf{Close}] \index{Range!Close} Touching distance: grapples, knife-work, hand on a relic.
\item[\textbf{Near}] \index{Range!Near} Same room/yard/deck; a rush away.
\item[\textbf{Far}] \index{Range!Far} Same site but distant; requires route or time to reach.
\item[\textbf{Absent}] \index{Range!Absent} Off-screen; requires scene change or significant effort to interact.
\end{description}

\subsection*{Movement Actions}
\label{sec:movement-actions}
\begin{itemize}
\item \textbf{Move}: Shift one range band as a \emph{beat}.
\item \textbf{Dash}: Shift two bands as your full action (terrain may require a roll).
\item \textbf{Melee Flag}: Mark when two parties are in Near range and directly engaged in combat.
\end{itemize}

\subsection*{Position States}
\index{Position!Dominant}\index{Position!Controlled}\index{Position!Desperate}
\begin{description}
\item[\textbf{Dominant}] \index{Position!Dominant} You have cover, leverage, or ritual footing. Failure still leaves options.
\item[\textbf{Controlled}] \index{Position!Controlled} Standard case: exposed lanes, rivals near, watchful eyes. Failure has teeth, but not ruin.
\item[\textbf{Desperate}] \index{Position!Desperate} Bad ground, bad odds, bad timing. Failure is severe; success may bring extra XP.
\end{description}

\paragraph{Position Shifting:}
\begin{itemize}
\item GM can spend \textbf{1 SB} to worsen Position by one step.
\item Player can spend \textbf{1 Boon} to improve Position by one step (once per action).
\item Narrative triggers (flanking, reinforcements, etc.) can shift Position without cost.
\end{itemize}

\section{Narrative Time}
\label{sec:narrative-time}
\index{Time!narrative}
Time is measured by \emph{importance} rather than duration.
\begin{description}
\item[\textbf{A Moment}] \index{Time!Moment} A glance, a strike, a whisper over a law-stone.
\item[\textbf{Some Time}] \index{Time!Some Time} A skirmish, a negotiation, a careful climb.
\item[\textbf{Significant Time}] \index{Time!Significant Time} Hours of march, rites, audits, stakeouts.
\item[\textbf{Days}] \index{Time!Days} Drills, recoveries, research, roadwork.
\end{description}

\section{Social Interactions}
\label{sec:social-interactions}
\index{Social}\index{Culture}
Social scenes use the same engine with \textbf{cultural color} (see Chapter \ref{ch:gm-core}).

\subsection*{Cultural Skill Emphases}
\label{sec:cultural-skills}
\index{Social!culture}
Different regions value different approaches. Use the table in \textbf{Chapter \ref{ch:lore-compendium} (World Regions and Cultures)} to see what skills are most effective where. In general:

\begin{center}
\begin{longtable}{ll}
\toprule
\textbf{Region} & \textbf{Valued Approach} \\
\midrule
Viterra & Rapport with parishes; Sway for markets; Command under writ. \\
Acasia & Rapport for kin-bridges; Command with banner-rights; Deceive risks honor timers. \\
Ecktoria & Sway in salons; Deceive at court; Perform in temple fora. \\
Kahfagia & Rapport aboard ship; Sway at piers; Command on a storming deck. \\
Vhasia & Presence for courtly etiquette; Sway for diplomacy; Insight for hidden loyalties. \\
Thepyrgos & Lore for academic debate; Sway for rhetorical argument; Investigation for uncovering truth. \\
Mistlands & Survival for reading fog signs; Resolve for resisting spirit influence; Ritual for bell-law. \\
Ykrul & Command for leadership; Performance for boast and saga; Survival for steppe navigation. \\
Sidhi & Sway for negotiation; Subterfuge for smuggling; Investigation for lighthouse justice. \\
Tulkani & Performance for storytelling; Insight for reading truth; Survival for road etiquette. \\
\bottomrule
\end{longtable}
\end{center}

\subsection*{Social Stakes \& Timers}
\label{sec:social-timers}
\index{Social!stakes}\index{Timers!social}
Your GM may use visible timers to track social progress:
\begin{itemize}
\item \textbf{Alliance Timer (Viterra):} Parishes and guilds come to your side.
\item \textbf{Honor Timer (Acasia):} Feasts, oaths, wyrd -- trust builds (or frays).
\item \textbf{Bureau Timer (Ecktoria):} Stamps, seals, approvals -- delay is pressure.
\item \textbf{Trust Timer (Kahfagia):} Pilots and crews extend favors and routes.
\item \textbf{Courtesy Timer (Aelinnel):} Fae negotiations -- speak true or lose standing.
\item \textbf{Hollow Attention Timer (Aelaerem/Mistlands):} The uninvited guest draws near.
\end{itemize}

\subsection*{Hospitality as a Social Mechanic}
\label{sec:hospitality}
\index{Social!hospitality}

In many cultures -- especially Aelaerem, Ykrul, Tulkani, Ubral, and Fhara -- hospitality is sacred. Accepting a cup, a meal, or a night's shelter creates a bond. Breaking hospitality is worse than breaking a sworn oath.

\paragraph{Hospitality Bond:} When you accept hospitality (food, drink, or shelter) from a host in a culture that honors guest-right, you and the host gain a mutual \textbf{Hospitality Bond [3]}. While the bond has segments, neither party may take hostile action against the other. The bond ticks when either party lies to the other, breaks a promise made under the host's roof, or acts to harm the other's interests. When it fills, the bond snaps -- and the violator gains a \textbf{Black Mark} in that culture (−1 die on all social rolls with that culture for one season).

\paragraph{Refusing Hospitality:} Refusing an offered cup, meal, or shelter is a grave insult in hospitality cultures. The GM may spend 1 SB to worsen Position by one step in all social interactions with that host and their community for the remainder of the visit.

\section{Supply and Resources}
\label{sec:supply}
\index{Supplies}\index{Supply Timer }

Track scarcity with a \textbf{Supply Timer } shared by the party's expedition.
\begin{center}
\begin{longtable}{cl}
\toprule
\textbf{Segments} & \textbf{State \& Effects} \\
\midrule
0 (Full) & Well-provisioned; no penalty. \\
2 (Low) & Minor frictions; --1 to resource checks. \\
3 (Dangerous) & Each PC gains \emph{Fatigue 1}. \\
4 (Empty) & Severe penalties; desperate measures. \\
\bottomrule
\end{longtable}
\end{center}

\subsection*{Foraging and Subsistence by Region}
\label{sec:foraging}
\index{Supplies!foraging}

Different regions offer different opportunities (and risks) for finding food and water.

\begin{center}
\begin{longtable}{lll}
\toprule
\textbf{Region} & \textbf{Forage DV} & \textbf{Typical Finds} \\
\midrule
Forest (Valewood, Aelinnel) & 2--3 & Berries, mushrooms, roots, small game \\
Steppe (Ykrul, Kuvani) & 3--4 & Edible roots, tubers, small game \\
Desert (Ashaan, Fhara) & 4--5 & Cactus water, insects, rare desert herbs \\
Marsh (Mistlands, Theona) & 3--4 & Cattails, frogs, medicinal mud \\
Mountain (Ubral, Pereshi) & 3--4 & Mountain herbs, small game, clean springs \\
Coast (Linn, Sidhi, Zakov) & 2--3 & Shellfish, seaweed, fish \\
Tundra (Northern Violet Steppe) & 4--5 & Lichens, frozen berries, small game \\
\bottomrule
\end{longtable}
\end{center}

\section{Engaging the World -- Player Actions}
\label{sec:player-actions-world}
\index{Procedures!world}
\begin{itemize}
\item \textbf{Scout \& Signal:} A follower can make the next travel action \emph{Dominant} (mark Exposure or Harm 1 on them -- see Chapter \ref{ch:running-assets-followers}).
\item \textbf{Local Color:} Briefly state what locals notice about you; GM offers a small fictional edge \emph{or} a tempting timer -- choose.
\item \textbf{Mark the Map:} On arrival, declare one change to the fiction (new ford, patron's shrine, toll-skip). GM may attach a minor timer as cost.
\item \textbf{Asset Activation:} Spend 1 Boon to activate an asset dramatically during a scene (see Chapter \ref{ch:running-assets-followers}).
\item \textbf{Follower Assistance:} Have a follower assist your actions for bonus dice (max +3 from all sources, see Chapter \ref{ch:gm-core}).
\item \textbf{Invoke Cultural Knowledge:} If your character has a relevant background or has spent significant time in a region, you may ask the GM one straightforward yes/no question about local customs, taboos, or safe practices.
\end{itemize}

\section{World Interaction -- Quick Reference}
\label{sec:world-interaction-quickref}

\begin{center}
\begin{longtable}{ll}
\toprule
\textbf{When you want to\dots} & \textbf{Use this approach} \\
\midrule
Travel quickly & Take the Guide role, prioritize Speed intent \\
Travel safely & Take the Scout role, prioritize Stealth or Survey intent \\
Learn local customs & Ask a local, invoke cultural knowledge, observe before acting \\
Gain a patron & Complete a favor, offer a service, impress with skill or honor \\
Avoid attention & Travel with Stealth intent, keep to back roads, pay bribes \\
Negotiate a dispute & Invoke common ground, offer a compromise, use social skills \\
Survive the wilds & Take the Quartermaster role, forage, find shelter \\
Honor hospitality & Accept offered food/drink, share a story, never draw first blood \\
Cross a thin place & Offer a memory, speak a true name, leave an offering \\
\bottomrule
\end{longtable}
\end{center}

\section{Summary}
\label{sec:world-summary}
\index{World interaction!summary}
The world has opinions. Movement is timers and color, position rises and sinks with weather and words, and every suit you draw speaks in a regional accent. Ask the land for a favor -- then pay it back on the road.

\paragraph{Remember:} Every interaction with the world is an opportunity. Use your assets, deploy your followers, and engage with the setting actively. The world responds to your choices, and every journey changes both you and the places you pass through. The road is not a ledger -- it is a conversation. Speak carefully. Listen closely. And always leave something behind for the next traveler.
